\exercisesection{}

\begin{enumerate}
\def\labelenumi{\arabic{enumi})}
\tightlist
\item
  Let \(\mathcal{T} = (\mathrm{K}_i, p_{ij})\) be an inverse system of compact spaces indexed by the set I, let K be a compact space, and let \((p_i)_{i \in I}\) be a coherent and separating family of continuous mappings \(p_i: \mathrm{K} \to \mathrm{K}_i\) .
\end{enumerate}

\begin{enumerate}
\def\labelenumi{\alph{enumi})}
\item
  Let A be the set of continuous functions on K of the form \(f_{i} \circ p_{i}\) , where i runs over I and \(f_{i}\) runs over the set of continuous real functions on \(K_{i}\) . Show that A is a dense linear subspace of \(\mathcal{C}(K)\) .
\item
  Let \((\mu_i)_{i \in I}\) be a sub-inverse system of positive measures on \(\mathcal{T}\) . Assume that the \(p_i\) are surjective. For every \(f \in A\) , let \(I_f\) be the set of \(i \in I\) such that \(f\) is of the form \(f_i \circ p_i\) with \(f_i \in \mathcal{C}(K_i)\) (necessarily unique). Show that there exists one and only one measure \(\pi\) on \(K\) such that \(\pi(f) = \inf_{i \in I_f} \mu_i(f_i)\) for all \(f \in A\) . Show that \(\pi\) is the
\end{enumerate}

greatest of the positive measures \(\mu\) on K such that \(p_i(\mu) \leqslant \mu_i\) for all \(i \in I\) .

¶ 2) Hypotheses and notations are those of Th. 1 of No.~2, for which we propose to indicate a new proof.

\begin{enumerate}
\def\labelenumi{\alph{enumi})}
\item
  Let K be a compact subset of T; for every \(i \in I\) , set \(\mathrm{K}_{i} = p_{i}(\mathrm{T})\) , and denote by \(q_{i}\) the mapping of K onto \(K_{i}\) that coincides on K with \(p_{i}\) ; for \(i \leqslant j\) , denote by \(q_{ij}\) the mapping of \(K_{j}\) into \(K_{i}\) that coincides on \(K_{j}\) with \(p_{ij}\) . Deduce from the preceding exercise the existence of a greatest positive measure \(\pi^{K}\) on K such that \(q_{i}(\pi^{K}) = (\mu_{i})_{\mathbf{K}_{i}}\) for all \(i \in I\) .
\item
  If K and L are compact subsets of T such that \(K \subset L\) , show that \((\pi^{\mathrm{L}})_{\mathrm{K}} \geqslant \pi^{\mathrm{K}}\) ; deduce from this that the set of measures of the form \(i^{\mathrm{K}}(\pi^{\mathrm{K}})\) (where K runs over the set of compact subsets of T, and \(i^{K}\) is the canonical injection of K into T) admits a supremum \(\mu\) in \(\mathcal{M}_{+}(\mathrm{T})\) .
\item
  Show that \(p_i(\mu) = \mu_i\) for every \(i \in I\) (it suffices to observe that \(p_i(\mu) \leqslant \mu_i\) and that the measures \(p_i(\mu)\) and \(\mu_i\) have the same total mass by the hypothesis (P)).
\end{enumerate}

